% ---------------------------------------------------------------------------
% Dataset and Exploratory Analysis (Section 3) - drafted for M3 (with M1's Phase 1 profiling).
% Numbers come from notebooks/01_eda.ipynb and results/metrics/; figures are
% results/figures/*.png (main.tex sets \graphicspath to ../results/).
% Needs in the preamble: \usepackage{booktabs, graphicx}.
% ---------------------------------------------------------------------------
\section{Dataset and Exploratory Analysis}
\label{sec:data}

\subsection{Data}
We use the Swahili Audio Classification dataset~\cite{zindi2021swahili}. It holds recordings of 12 spoken Swahili
words, made by about 300 crowd-workers in Kenya: the digits \textit{moja} to \textit{kumi} (one to ten),
\textit{ndio} (yes) and \textit{hapana} (no). There are 4{,}200 labelled clips, exactly 350 per word, and
1{,}800 unlabelled competition test clips. All 6{,}000 files are 16~kHz mono 16-bit PCM, with no corrupt files
and no byte-identical duplicates. The competition allows no external data and forbids redistribution, so the audio
is kept off the public repository. As described in Section~\ref{sec:method}, we froze one stratified 70/15/15
split (2{,}940 / 630 / 630 clips). The distribution statistics below use the training split only. Analyses that
score a model (the confusions, the prefix curve and the normalisation check) use the validation split, as model
selection does. None uses the test split.

\subsection{What the data revealed}
\paragraph{The words are short, but the clips are long}
Clips last 2.5--62~s (median 4.3~s), while the spoken word lasts a median 0.55~s. Word duration ranges from
0.45~s for \textit{tatu} to 0.65~s for the three-syllable \textit{hapana}, and 89.5\% of words last at most 1.0~s
(Fig.~\ref{fig:durations}). A 30~dB silence trim leaves a median of 1.9~s and a 95th percentile of 5.5~s,
because background noise is kept as part of the ``word''. The models therefore read a centred window around the
clip's highest-energy region: 1.5~s by default, 1.0~s for the final A3 and 2.0~s for the final A4. The 1.5~s window
retains a median 99.8\% of each clip's above-noise energy.

\paragraph{Recordings vary in level, but are mostly clean}
Levels span about 19~dB between the 5th and 95th percentiles (median $-32$~dBFS), and 3.8\% of clips reach full
scale, so windows are peak-normalised. The median estimated SNR is 42~dB, and 80\% of clips exceed 30~dB. Our
synthetic noise augmentation (10--30~dB) is therefore harsher than most real recordings. The noisy minority matters,
however: the models fail most on clips below 30~dB (Section~\ref{sec:errors}).

\paragraph{The vocabulary contains an anagram}
\textit{sita} (six) and \textit{tisa} (nine) contain the same sounds in a different order. The average
spectrograms show it directly (Fig.~\ref{fig:spectra}): the high-frequency energy of /s/ comes at the start of
\textit{sita} and in the middle of \textit{tisa}. Other near-pairs are \textit{nne}/\textit{nane} (letter edit
distance 1), and \textit{tano}/\textit{tatu} and \textit{saba}/\textit{sita} (distance 2). We track these four
pairs for every model. Spelling similarity, however, only weakly predicts what an order-free model confuses
(Spearman $\rho=-0.18$ over all 66 pairs, $p=0.16$). Its most confused pair on validation, \textit{mbili}/\textit{ndio}, is far
apart in spelling, but both words begin with a prenasalised stop and contain /i/.

\paragraph{Order-free features do not separate the words}
A t-SNE map of order-free MFCC statistics shows no word forming its own cluster. The only distinct islands contain
clips of many words, so they probably group recordings by speaker or recording condition. This is consistent with
the order-free baseline's ceiling of 77.9\%, and with the large gain from per-utterance normalisation in the HMM
(53\% to 76\% accuracy).

\paragraph{The evidence builds up through the word}
Trained on only the first 25/50/75/100\% of each word, the order-free classifier reaches 49/72/80/80\% accuracy.
Words that share a first syllable (\textit{tano}, \textit{tatu}) separate only once their second syllable begins.
The evidence therefore builds up over most of the word, not just its onset, and the window must cover the whole
word.

\paragraph{No leakage, but speakers are unknown}
No pair of clips across the splits is a near-duplicate (the maximum cosine similarity of standardised MFCC
statistics is 0.94). Median pitch spans 106--235~Hz between the 10th and 90th percentiles, consistent with a mix of
adult male and female voices. Speaker identities are not provided, however, so the same speaker may appear in
both training and test clips (Section~\ref{sec:limitations}).

\begin{table}[t]
  \centering
  \caption{How the exploratory analysis shaped the design.}
  \label{tab:eda}
  \small
  \begin{tabular}{@{}p{0.47\linewidth}p{0.45\linewidth}@{}}
    \toprule
    Finding & Decision \\
    \midrule
    350 clips per word, perfectly balanced & No class weighting; accuracy, macro-F1 and log loss \\
    Words last a median 0.55~s in 2.5--62~s clips; silence trimming leaves noisy tails & Centred 1.5~s energy window; 1.0/2.0~s tested \\
    Levels vary by $\sim$19~dB; median SNR 42~dB & Peak normalisation; CMVN for the HMM \\
    \textit{sita}/\textit{tisa} differ only in where /s/ occurs & Sequence models; four sound-alike pairs tracked \\
    Order-free features do not cluster by word & Per-utterance normalisation; order-aware models \\
    Evidence arrives through the word & Models see the whole word \\
    No near-duplicates; no speaker IDs & No leakage found; speaker overlap is a limitation \\
    \bottomrule
  \end{tabular}
\end{table}

\begin{figure}[t]
  \centering
  \includegraphics[width=\linewidth]{figures/eda_durations.png}
  \caption{Left: whole clips are far longer than the spoken words they contain (training split, log scale).
  Right: estimated word duration per word.}
  \label{fig:durations}
\end{figure}

\begin{figure}[t]
  \centering
  \includegraphics[width=\linewidth]{figures/eda_mean_spectrograms.png}
  \caption{Average log-mel spectrogram of each word in the centred 1.5~s window. The high-frequency /s/ comes at the
  start of \textit{sita} and in the middle of \textit{tisa}.}
  \label{fig:spectra}
\end{figure}
